\section{源码等价的时序优化与生命周期模型}
\label{sec:time-series-source-equivalent}

本章只描述 \srcpath{src/time\_series/} 当前执行的模型。主调度模型由
\texttt{build\_uc\_milp} 构造；它不是单一的教材 SCUC，而是由选项门控的线性混合整数模型。
令 $T=\texttt{num\_steps}$、$\Delta t=\texttt{step\_duration\_hr}$。当 $T>0$ 时，代码要求
$\Delta t$ 为有限正数；配置含重复 profile id 或非有限 profile 值时直接抛出异常。依据：
\srcpath{src/time\_series/time\_series\_pf.cpp:find\_profile}。

\subsection{模型启用条件与变量块}
代码中的三个主要门控谓词为
\begin{align}
 use_{acnet}&=\texttt{enable\_network\_constraints}\land(n_{acbus}>1),\\
 use_{dcnet}&=use_{acnet}\land\texttt{enable\_dc\_network\_constraints}
 \land(n_{dcbus}>0)\land(C>0),\\
 use_{shift}&=\texttt{enable\_demand\_response}\land(F>0)\land
 \texttt{dr\_shiftable}.
\end{align}
依据：\srcpath{src/time\_series/time\_series\_pf.cpp:build\_uc\_milp}。

变量按源码固定顺序拼接。基础块为发电出力 $p_g$、二进制开机 $u_g$、启动费用变量 $s_g$、
AC/DC 储能净出力及期末 SOC、可再生出力；可选块包括外部电网、需求响应升降量、可再生削减、
静态电源、微网交换与状态、DC 运输流、储能吞吐、VPP、电能路由器、移动储能、充电量与模式、
非平衡母线相角、VSC 和 DC--DC 功率。每块维数及偏移严格见
\srcpath{src/time\_series/time\_series\_pf.cpp:MobStor}，本章不改变这一索引空间。

\subsection{目标函数的实际权重}
模型始终最小化 $\mathbf c^T\mathbf x$。目标模式只按下式设置四个权重：
\begin{equation}
 (w_c,w_e,w_l,w_r)=
 \begin{cases}
 (1,0,0,0),&\texttt{Cost},\\
 (0,1,0,1),&\texttt{Carbon},\\
 (10^{-3},0,0,1),&\texttt{MinCurtailment},\\
 (0,0,1,1),&\texttt{MinLoss},\\
 (\texttt{w\_cost},\texttt{w\_carbon},\texttt{w\_loss},
  \texttt{w\_curtailment}),&\texttt{Weighted}.
 \end{cases}
\end{equation}
对常规机组，逐时系数精确为
\begin{align}
 c(p_{g,t})&=(w_c c_{1,g}+w_e e_g+w_l)\Delta t,\\
 c(u_{g,t})&=w_c\max(0,c_{0,g})\Delta t,\\
 c(s_{g,t})&=w_c.
\end{align}
固定 DC 电源、PV 和未调度静态电源的费用进入 \texttt{obj\_offset}，因此求解器目标并不等于
最终报告目标；报告值加回该常数。依据：
\srcpath{src/time\_series/time\_series\_pf.cpp:dcdc\_idx（lambda）}。

可再生机组变量范围与系数为
\begin{equation}
 0\le p_{r,t}\le profile_{r,t}P_r^{rated},\qquad
 c(p_{r,t})=[w_c c_{1,r}-(w_c c^{curt}_r+w_r)]\Delta t.
\end{equation}
外部电网变量范围为 $[-\max(1,\texttt{external\_grid\_cap\_mw}),
+\max(1,\texttt{external\_grid\_cap\_mw})]$，价格 profile 存在时直接使用该绝对价格，否则用
\texttt{cost\_c1}。需求响应的上下调容量分别为
\begin{equation}
 \bar p^{up}=\max(0,flex^{up})a,\qquad
 \bar p^{dn}=\min(\max(0,flex^{dn})a,\max(0,p_0)),\quad
 a=\operatorname{clamp}(availability/100,0,1),
\end{equation}
且两者系数均为 $\max(0,\texttt{w\_demand\_response})\Delta t$。依据：
\srcpath{src/time\_series/time\_series\_pf.cpp:dcdc\_idx（lambda）}。

\subsection{机组约束}
出力边界由两行线性不等式实现：
\begin{equation}
 p_{g,t}-P_g^{max}u_{g,t}\le0,\qquad
 -p_{g,t}+P_g^{min}u_{g,t}\le0.
\end{equation}
爬坡上、下限分别为
\begin{align}
 p_{g,t}-p_{g,t-1}&\le R_g^{up}\Delta t\,60,\\
 p_{g,t-1}-p_{g,t}&\le R_g^{dn}\Delta t\,60.
\end{align}
若计算出的任一爬坡限值非正，代码将其替换为 $P_g^{max}$，并非删除该约束。启动费用变量满足
\begin{align}
 -s_{g,0}+C_g^{start}u_{g,0}&\le C_g^{start}u_g^{init},
 \quad u_g^{init}=\mathbb 1[p_g^{initial}>10^{-6}],\\
 -s_{g,t}+C_g^{start}(u_{g,t}-u_{g,t-1})&\le0\quad(t\ge1).
\end{align}
依据：\srcpath{src/time\_series/time\_series\_pf.cpp:stor\_power\_idx（lambda）}。

令 $N_g^{up}=\lceil T_g^{up}/\Delta t\rceil$。仅当 $N_g^{up}\ge2$ 时，对
$t\ge1$ 和 $1\le k\le\min(N_g^{up}-1,T-1-t)$ 加入
\begin{equation}
 u_{g,t}-u_{g,t-1}-u_{g,t+k}\le0.
\end{equation}
最小停机时间同理，以
\begin{equation}
 u_{g,t-1}-u_{g,t}+u_{g,t+k}\le1
\end{equation}
实现。依据：\srcpath{src/time\_series/time\_series\_pf.cpp:stor\_power\_idx（lambda）}。

\subsection{储能净出力拆分与 SOC}
代码定义有符号净出力 $p=p^{dis}-p^{ch}$，其中 $p^{ch}\ge0$；实际线性约束为
\begin{align}
 -p-p^{ch}&\le0,\\
 p+p^{ch}-\bar P^{dis}z&\le0,\\
 p^{ch}+\bar P^{ch}z&\le\bar P^{ch}.
\end{align}
只有充、放两个方向容量均大于 $10^{-9}$，且任一效率与 $1$ 相差超过 $10^{-12}$ 时，$z$ 才是
二进制变量；否则它是 $[0,1]$ 连续变量。依据：
\srcpath{src/time\_series/time\_series\_pf.cpp:dcdc\_idx（lambda）}、
\srcpath{src/time\_series/time\_series\_pf.cpp:stor\_power\_idx（lambda）}。

自放电保持率不是线性近似，而是
\begin{equation}
 \rho=\left[\operatorname{clamp}(1-self\_discharge\_pct/100,0,1)\right]^
 {\max(0,\Delta t)}.
\end{equation}
令 $E'=E$ 当 $E\ge10^{-12}$，否则 $E'=1$；
$\eta_c,\eta_d$ 分别截断到 $[10^{-6},1]$。AC 与 DC 储能均使用同一行：
\begin{equation}
 SOC_t+\frac{\Delta t}{\eta_dE'}p_t+
 \frac{\Delta t}{E'}(\eta_d^{-1}-\eta_c)p_t^{ch}
 -\rho SOC_{t-1}=0,
\end{equation}
其中 $t=0$ 时删除上一状态项并令右端为 $\rho SOC^{init}$。若启用循环边界，最后一个
SOC 变量的上下界同时固定为
$\operatorname{clamp}(SOC^{init},SOC^{min},SOC^{max})$。依据：
\srcpath{src/time\_series/time\_series\_pf.cpp:storage\_retention}、
\srcpath{src/time\_series/time\_series\_pf.cpp:dcdc\_idx（lambda）}、
\srcpath{src/time\_series/time\_series\_pf.cpp:stor\_power\_idx（lambda）}。

吞吐辅助量只在退化或储能报价需要时创建，并满足
\begin{equation}
 p_{k,t}-g_{k,t}\le0,\qquad-p_{k,t}-g_{k,t}\le0.
\end{equation}
启用退化且 \texttt{daily\_cycle\_limit}$>0$ 时，再加入
\begin{equation}
 \sum_tg_{k,t}\Delta t\le2\,daily\_cycle\_limit_k\,\tilde E_k,
\end{equation}
其中 $\tilde E_k=E_k$ 当 $E_k>10^{-9}$，否则 $\tilde E_k=1$；此守卫阈值为 $10^{-9}$，
与 SOC 式中 $E'$ 的 $10^{-12}$ 不同。
依据：\srcpath{src/time\_series/time\_series\_pf.cpp:stor\_power\_idx（lambda）}。

\subsection{功率平衡与网络约束边界}
当 $use_{acnet}=false$ 时，代码使用一条铜板平衡。其左端按源码符号累加：常规机组、AC/DC 储能、
可再生、外部电网、需求响应下调、静态电源、微网出口、VPP、路由器净输出和已连接移动储能为正；
需求响应上调、可再生削减和路由器输入为负；右端为 \texttt{total\_load[t]}。完整项见
\srcpath{src/time\_series/time\_series\_pf.cpp:stor\_power\_idx（lambda）}。

当 $use_{acnet}=true$ 时，逐 AC 母线装配同类平衡，并使用非平衡母线相角。投运、额定值为正、
$|x|\ge10^{-12}$ 且两端有效的 AC 支路才加入限额：
\begin{equation}
 -S_{ij}^{max}\le\frac{S_{base}}{x_{ij}}(\theta_i-\theta_j)
 \le S_{ij}^{max}.
\end{equation}
依据：\srcpath{src/time\_series/time\_series\_pf.cpp:stor\_power\_idx（lambda）}、
\srcpath{src/time\_series/time\_series\_pf.cpp:stor\_power\_idx（lambda）}。

DC 网络模型没有 DC 电压变量。启用显式 DC 运输流时，每条合格支路只创建
$f_{ij,t}\in[-rate_{ij},rate_{ij}]$ 并进入节点平衡；否则代码明确令
\texttt{n\_ineq\_dc\_line}=0，依赖换流器变量边界，不能把它写成电阻潮流或独立 DC 支路热约束。
依据：\srcpath{src/time\_series/time\_series\_pf.cpp:SgenSupply}、
\srcpath{src/time\_series/time\_series\_pf.cpp:stor\_power\_idx（lambda）}、
\srcpath{src/time\_series/time\_series\_pf.cpp:stor\_power\_idx（lambda）}。

\subsection{UC 到逐步 OPF 的真实跟踪带}
仅在应用 UC 日程、机组在役且 \texttt{enable\_uc\_opf\_tracking\_band} 为真时修改 OPF 边界。
对非松弛机组可切换到专用相对/绝对参数；实际带宽是
\begin{equation}
 b_g=\max\left(abs_g,rel_g\max(|p_g^{UC}|,1)\right),
\end{equation}
并赋值
\begin{equation}
 P_g^{min,new}=\max(P_g^{min},p_g^{UC}-b_g),\qquad
 P_g^{max,new}=\min(P_g^{max},p_g^{UC}+b_g).
\end{equation}
若新上界小于新下界，两者都固定到
$\operatorname{clamp}(p_g^{UC},P_g^{min},P_g^{max})$。依据：
\srcpath{src/time\_series/time\_series\_pf.cpp:build\_time\_series\_system\_snapshot}。

\subsection{按日并行的实际配置}
日模式只按下列赋值切换：\texttt{SCUC} 为
\texttt{skip\_uc=false, run\_opf=true, fix\_commitment=false}；
\texttt{DynamicSCED} 为 \texttt{false,false,true}；
\texttt{DynamicOPF} 为 \texttt{true,true,false}。并行模式强制 AC OPF 后端为
\texttt{ParityIPM}。每日末 SOC 是否固定，完全由
\texttt{enforce\_daily\_cyclic\_soc} 传给
\texttt{enforce\_terminal\_soc\_cyclic}，不是无条件成立。依据：
\srcpath{src/time\_series/annual\_production\_sim.cpp:slice\_2d\_int（lambda）}。

\subsection{生命周期赋值式}
生命周期循环的年序号从 $y=1$ 开始。负荷和 PV 使用
\begin{align}
 factor_y^{load}&=(1+g)^{y-1},\\
 factor_y^{pv}&=(1-d_{pv})^{y-1},\\
 p_{load,y}&=p_{load,0}factor_y^{load},\quad
 q_{load,y}=q_{load,0}factor_y^{load},\\
 P_{pv,y}^{max}&=P_{pv,0}^{max}factor_y^{pv},\quad
 p_{pv,y}=\min(p_{pv,y},P_{pv,y}^{max}).
\end{align}
依据：\srcpath{src/time\_series/lifecycle\_simulation.cpp:OrigStorage}。

AC 储能的两个健康度分量及合成值是
\begin{align}
 SOH_y^{cal}&=\max(0,1-d_{cal}(y-1)),\\
 SOH_y^{cyc}&=\max(0,1-d_{cyc}N_{cycles}),\\
 SOH_y&=\min(SOH_y^{cal},SOH_y^{cyc}).
\end{align}
EOL 阈值先做百分数归一化：$EOL'=EOL/100$ 当 $EOL>1$（视为百分数），否则 $EOL'=EOL$；
当 $SOH_y\le\operatorname{clamp}(EOL',0,1)$ 时，代码立即记录更换并将两个分量、合成值置为 $1$，
循环数置零；随后令 $E_y^{rated}=E_0^{rated}SOH_y$、$E_y=E_y^{rated}SOC^{init}$。
更换成本精确为 \texttt{replacement\_cost}$\times E_0^{rated}\times1000$。依据：
\srcpath{src/time\_series/lifecycle\_simulation.cpp:OrigStorage}。

年度成本和当年更换成本使用同一折现因子
\begin{equation}
 d_y=(1+r)^{-(y-1)},\qquad
 NPV\leftarrow NPV+d_y\left(C_y+\sum_{e\in replacements_y}C_e\right).
\end{equation}
碳量只累加，不折现。依据：
\srcpath{src/time\_series/lifecycle\_simulation.cpp:OrigStorage}。

\subsection{覆盖边界}
本章逐式覆盖主 MILP 的核心变量、目标、机组约束、储能与网络门控、OPF 跟踪带、按日模式和
生命周期状态更新。能量路由器、VPP、移动储能、微网、需求响应等可选块的完整逐行矩阵仍以
\srcpath{src/time\_series/time\_series\_pf.cpp:stor\_power\_idx（lambda）} 与
\srcpath{src/time\_series/time\_series\_pf.cpp:stor\_power\_idx（lambda）} 为准；本章没有用通用模型补写未展开的行。
